\documentclass[a4paper,11pt]{article}
\usepackage[utf8]{inputenc}
\usepackage[T1]{fontenc}
\usepackage[spanish]{babel}
\usepackage{geometry}
\usepackage{array}
\usepackage{colortbl}
\usepackage{xcolor}
\usepackage{booktabs}
\usepackage{enumitem}
\usepackage{fancyhdr}
\usepackage{tikz}

\geometry{left=2cm, right=2cm, top=2.5cm, bottom=2cm}

\definecolor{violeta}{HTML}{7B2D8E}
\definecolor{azulg}{HTML}{3498DB}
\definecolor{rojog}{HTML}{E74C3C}
\definecolor{gris}{HTML}{F2F2F2}
\definecolor{verde}{HTML}{27AE60}
\definecolor{naranja}{HTML}{E67E22}

\pagestyle{fancy}
\fancyhf{}
\fancyhead[L]{\small\textbf{Nuevo Compañeros 1 — U03: La Familia}}
\fancyhead[R]{\small\textbf{Píldora formativa 3.7}}
\fancyfoot[C]{\thepage}

\setlength{\parindent}{0pt}

\begin{document}

\begin{center}
{\LARGE\textbf{PÍLDORA FORMATIVA 3.7}}\\[0.3cm]
{\Large Fonética: el sonido /\texttheta/}\\[0.2cm]
{\normalsize Sección Comunicación (p.39) — Proyectable — 4 diapositivas}\\[0.2cm]
{\small Versión enriquecida secuencial para Fases P1--P5}
\end{center}

\vspace{0.5cm}

\noindent\fbox{\parbox{\dimexpr\linewidth-2\fboxsep-2\fboxrule}{%
\textbf{Contenido para el profesor}\\[0.2cm]
El sonido /\texttheta/ (interdental fricativo sordo) se produce colocando la punta de la lengua entre los dientes superiores e inferiores y expulsando el aire suavemente. Es el mismo sonido que el inglés \textit{th} en \textit{think, thick, three} (no el de \textit{this, that}, que es sonoro /\dh/).\\[0.2cm]
\textbf{Distribución ortográfica:}\\
/\texttheta/ + a, o, u $\rightarrow$ se escribe con \textbf{Z}: za, zo, zu (zapato, azul, zumo)\\
/\texttheta/ + e, i $\rightarrow$ se escribe con \textbf{C}: ce, ci (cine, cinco, bicicleta)\\[0.2cm]
\textbf{Nota:} Esta regla aplica al español peninsular estándar. En Hispanoamérica y sur de España: seseo (/s/).
}}

\vspace{0.8cm}

%% ============================================================
%% DIAPOSITIVA 1
%% ============================================================

\noindent\textcolor{violeta}{\rule{\linewidth}{2.5pt}}

\section*{DIAPOSITIVA 1 — ESCUCHAD ESTE SONIDO}

{\small Fase P1 (Escuchar modelo) — Técnica: Input auditivo saturado con articulación visible}\\
{\small\textit{Principio:} El estudiante recibe 5 exposiciones al sonido /\texttheta/ en palabras reales, con énfasis articulatorio visible. Las tarjetas con las palabras escritas y la letra resaltada en color conectan el sonido con la grafía desde el primer momento.}

\vspace{0.5cm}

\subsection*{En pantalla}

5 tarjetas de palabras en horizontal:

\begin{center}
{\Large
a\textbf{\textcolor{rojog}{z}}ul \quad\quad
bi\textbf{\textcolor{rojog}{c}}i\textbf{\textcolor{rojog}{c}}leta \quad\quad
\textbf{\textcolor{rojog}{c}}in\textbf{\textcolor{rojog}{c}}o \quad\quad
mar\textbf{\textcolor{rojog}{z}}o \quad\quad
\textbf{\textcolor{rojog}{z}}apato
}
\end{center}

\vspace{0.3cm}

Las letras \textbf{z} y \textbf{c} que representan /\texttheta/ están en rojo y en tamaño grande.

\vspace{0.3cm}

Debajo: icono de oreja + \textit{Escuchad y mirad la boca del profesor.}

\vspace{0.2cm}

Símbolo fonético: {\Large /\texttheta/} (entre barras, tamaño grande)

\vspace{0.2cm}

Instrucción articulatoria: \textit{Lengua entre los dientes. Aire sale suavemente.}

\vspace{0.5cm}

\subsection*{Instrucciones para el profesor}

\begin{enumerate}[leftmargin=1.5cm, itemsep=4pt]
\item Diga cada palabra lentamente, exagerando el sonido /\texttheta/. Muestre la posición de la lengua (entre los dientes).
\item Repita las 5 palabras dos veces.
\item Ponga el audio (pista 39). Los estudiantes escuchan y observan las tarjetas.
\item Pida: \textit{¿Qué sonido se repite en todas estas palabras?} Confirme: \textit{Sí, es el sonido /\texttheta/.} Produzca el sonido aislado. Pida que intenten producirlo.
\item Si dispone de un espejo, ofrezca que se miren la lengua al producir el sonido.
\end{enumerate}

\newpage

%% ============================================================
%% DIAPOSITIVA 2
%% ============================================================

\noindent\textcolor{violeta}{\rule{\linewidth}{2.5pt}}

\section*{DIAPOSITIVA 2 — ¿CUÁL OYES?}

{\small Fase P2 (Discriminar) — Técnica: Pares mínimos receptivos + descubrimiento inductivo de regla ortográfica}\\
{\small\textit{Principio:} El estudiante ve dos palabras que difieren en un solo sonido (/s/ vs /\texttheta/) y debe identificar cuál ha dicho el profesor. Después, las palabras con /\texttheta/ se clasifican por letra (z o c) y el estudiante descubre la regla ortográfica de forma inductiva.}

\vspace{0.5cm}

\subsection*{En pantalla — Parte 1: Discriminación}

8 pares mínimos /s/ vs /\texttheta/. El profesor dice UNA palabra de cada par:

\begin{center}
\renewcommand{\arraystretch}{1.4}
\begin{tabular}{|c|>{\centering\arraybackslash}p{4cm}|>{\centering\arraybackslash}p{4cm}|}
\hline
\rowcolor{gris}
 & \textbf{Opción A (/s/)} & \textbf{Opción B (/\texttheta/)} \\
\hline
1 & ca\textbf{s}a & ca\textbf{\textcolor{rojog}{z}}a \\
\hline
2 & a\textbf{s}a & ha\textbf{\textcolor{rojog}{z}}a \\
\hline
3 & a\textbf{s}ado & a\textbf{\textcolor{rojog}{z}}ada \\
\hline
4 & \textbf{s}ien & \textbf{\textcolor{rojog}{c}}ien \\
\hline
5 & \textbf{s}eta & \textbf{\textcolor{rojog}{z}}eta \\
\hline
6 & co\textbf{s}er & co\textbf{\textcolor{rojog}{c}}er \\
\hline
7 & \textbf{s}umo & \textbf{\textcolor{rojog}{z}}umo \\
\hline
8 & \textbf{s}errar & \textbf{\textcolor{rojog}{c}}errar \\
\hline
\end{tabular}
\end{center}

\vspace{0.5cm}

\subsection*{En pantalla — Parte 2: Clasificación}

Las palabras con /\texttheta/ se clasifican en dos columnas:

\begin{center}
\renewcommand{\arraystretch}{1.4}
\begin{tabular}{|>{\centering\arraybackslash}p{5cm}|>{\centering\arraybackslash}p{5cm}|}
\hline
\rowcolor{naranja!15}
\textbf{Con Z} & \textbf{Con C} \\
\hline
ca\textbf{z}a & \textbf{c}ien \\
ha\textbf{z}a & \textbf{c}errar \\
a\textbf{z}ada & co\textbf{c}er \\
\textbf{z}eta & \\
\textbf{z}umo & \\
\hline
\end{tabular}
\end{center}

\vspace{0.3cm}

Pregunta inductiva: \textit{¿Qué vocales vienen después de la Z?} $\rightarrow$ a, o, u.\\
\textit{¿Qué vocales vienen después de la C?} $\rightarrow$ e, i.

\vspace{0.3cm}

Tras clic, aparece la regla:
\begin{itemize}[leftmargin=2cm]
\item /\texttheta/ + \textbf{a, o, u} $\rightarrow$ \textbf{Z}: \textbf{z}a, \textbf{z}o, \textbf{z}u
\item /\texttheta/ + \textbf{e, i} $\rightarrow$ \textbf{C}: \textbf{c}e, \textbf{c}i
\end{itemize}

\vspace{0.5cm}

\subsection*{Instrucciones para el profesor}

\begin{enumerate}[leftmargin=1.5cm, itemsep=4pt]
\item Parte 1: Muestre cada tarjeta con dos palabras e imágenes. Diga UNA de las dos. Los estudiantes señalan cuál han oído. Varíe el orden.
\item Parte 2: Separe las palabras con /\texttheta/ y pida que las clasifiquen en dos columnas (Z / C). Haga las preguntas inductivas. Deje que descubran la regla.
\item Confirme la regla en la pizarra: /\texttheta/ + a,o,u = Z; /\texttheta/ + e,i = C.
\item Añada: \textit{¿Existe 'ze' o 'zi' en español?} Respuesta: casi nunca. \textit{¿Y 'ca'?} Sí, pero suena /ka/, no /\texttheta a/.
\end{enumerate}

\newpage

%% ============================================================
%% DIAPOSITIVA 3
%% ============================================================

\noindent\textcolor{violeta}{\rule{\linewidth}{2.5pt}}

\section*{DIAPOSITIVA 3 — REPETID CONMIGO}

{\small Fase P3 (Repetir con apoyo) — Técnica: Repetición coral $\rightarrow$ individual $\rightarrow$ teléfono fonético $\rightarrow$ velocidad creciente}\\
{\small\textit{Principio:} La repetición coral reduce la ansiedad. La repetición individual permite feedback personalizado. El teléfono fonético gamifica la producción del sonido. La velocidad creciente automatiza la articulación.}

\vspace{0.5cm}

\subsection*{En pantalla}

Las 5 palabras modelo en tamaño grande, con /\texttheta/ resaltado:

\begin{center}
{\Large
a\textbf{\textcolor{rojog}{Z}}ul \quad\quad
bi\textbf{\textcolor{rojog}{C}}i\textbf{\textcolor{rojog}{C}}leta \quad\quad
\textbf{\textcolor{rojog}{C}}in\textbf{\textcolor{rojog}{C}}o \quad\quad
mar\textbf{\textcolor{rojog}{Z}}o \quad\quad
\textbf{\textcolor{rojog}{Z}}apato
}
\end{center}

\vspace{0.5cm}

Cuatro fases indicadas con iconos:

\begin{center}
\renewcommand{\arraystretch}{1.8}
\begin{tabular}{|>{\centering\arraybackslash}p{3cm}|>{\raggedright\arraybackslash}p{9cm}|}
\hline
\rowcolor{verde!15}
\textbf{Fase A} & TODOS JUNTOS (coral) \\
\hline
\textbf{Fase B} & UNO A UNO (individual) \\
\hline
\rowcolor{naranja!10}
\textbf{Fase C} & ¡TELÉFONO FONÉTICO! (cadena en susurro) \\
\hline
\textbf{Fase D} & ¡MÁS RÁPIDO! (velocidad creciente) \\
\hline
\end{tabular}
\end{center}

\vspace{0.5cm}

\subsection*{Instrucciones para el profesor}

\begin{enumerate}[leftmargin=1.5cm, itemsep=4pt]
\item \textbf{Coral} (30 seg): Diga cada palabra. Toda la clase repite al unísono. Repita dos veces la serie completa.
\item \textbf{Individual} (1 min): Señale a estudiantes individuales. Si la pronunciación no es correcta, modele exagerando: \textit{Mirad: la lengua entre los dientes. /\texttheta/apato. Intentadlo.} No corrija más de una vez por estudiante.
\item \textbf{Teléfono fonético} (1 min): Organice filas de 4--5 estudiantes. Susurre una palabra al primero. Cada estudiante la susurra al siguiente. El último la dice en voz alta. Si la palabra llega correcta, la fila gana un punto. Haga 3 rondas (zapato, bicicleta, cinco).
\item \textbf{Velocidad creciente} (30 seg): Diga las 5 palabras cada vez más rápido. Los estudiantes siguen el ritmo. Termine a velocidad normal.
\end{enumerate}

\newpage

%% ============================================================
%% DIAPOSITIVA 4
%% ============================================================

\noindent\textcolor{violeta}{\rule{\linewidth}{2.5pt}}

\section*{DIAPOSITIVA 4 — COMPLETAD CON Z O C}

{\small Fase P4--P5 (Producir en palabras nuevas + Ortografía asociada) — Técnica: Producción en contexto nuevo + regla ortográfica con weaning off}\\
{\small\textit{Principio:} El estudiante aplica el sonido /\texttheta/ en 10 palabras nuevas. Primero completa los huecos con la regla visible, después sin la regla. El audio de verificación (pista 40) cierra el ciclo.}

\vspace{0.5cm}

\subsection*{En pantalla}

10 palabras con hueco para z o c (de la Actividad 9):

\begin{center}
\renewcommand{\arraystretch}{1.5}
\begin{tabular}{|c|>{\centering\arraybackslash}p{4cm}|>{\centering\arraybackslash}p{4cm}|}
\hline
\rowcolor{gris}
 & \textbf{Palabra con hueco} & \textbf{Respuesta} \\
\hline
1 & \_ine & \textbf{c}ine \\
\hline
2 & on\_e & on\textbf{c}e \\
\hline
3 & \_umo & \textbf{z}umo \\
\hline
4 & \_oo & \textbf{z}oo \\
\hline
5 & balon\_esto & balon\textbf{c}esto \\
\hline
6 & \_iclismo & \textbf{c}iclismo \\
\hline
7 & \_iudad & \textbf{c}iudad \\
\hline
8 & \_ona & \textbf{z}ona \\
\hline
9 & pi\_arra & pi\textbf{z}arra \\
\hline
10 & habita\_ión & habita\textbf{c}ión \\
\hline
\end{tabular}
\end{center}

\vspace{0.3cm}

\noindent\colorbox{verde!15}{\parbox{\dimexpr\linewidth-2\fboxsep}{%
\centering\textbf{Regla:} /\texttheta/ + a, o, u = \textbf{Z} \quad|\quad /\texttheta/ + e, i = \textbf{C}
}}

\vspace{0.3cm}

Botón: \textit{Ocultar regla} (para weaning off)

\vspace{0.5cm}

\subsection*{Instrucciones para el profesor}

\begin{enumerate}[leftmargin=1.5cm, itemsep=4pt]
\item \textbf{Con regla visible:} Pida que completen los huecos individualmente (1.5 min). Pueden consultar la regla.
\item \textbf{Weaning off:} Oculte la regla. Lea las 3 últimas palabras en voz alta y pida que digan si llevan z o c sin mirar la regla.
\item Ponga el audio (pista 40) para verificar. Los estudiantes escuchan y corrigen.
\item Puesta en común: lea cada palabra y pida que la repitan pronunciando /\texttheta/ correctamente.
\item \textbf{Microtarea oral en parejas:} cada estudiante elige 3 palabras y dice una frase con cada una (ej: \textit{Yo vivo en una ciudad}, \textit{Tengo once años}). El compañero escucha y levanta el pulgar si pronuncia /\texttheta/ correctamente.
\item \textbf{Dictado de transferencia:} diga 5 palabras nuevas: \textit{cereza, azúcar, cero, zarpa, circo}. Los estudiantes escriben. Comprueban con la regla. Diga: \textit{¿Veis? La regla funciona siempre.}
\end{enumerate}

\end{document}
